% ============================================================
\chapter{Optimisation de Portefeuille}
\label{ch:portefeuille}
% ============================================================

En 1952, Harry Markowitz, alors \'etudiant en doctorat \`a l'universit\'e de Chicago, publie un article de quatorze pages qui va r\'evolutionner la finance~: \og Portfolio Selection \fg. Son id\'ee est simple mais profonde~: un investisseur ne devrait pas choisir ses actifs individuellement, mais optimiser son \emph{portefeuille} dans son ensemble, en tenant compte des corr\'elations entre les rendements. La diversification --- ``ne pas mettre tous ses \oe ufs dans le m\^eme panier'' --- re\c{c}oit ainsi une fondation math\'ematique rigoureuse, exprim\'ee comme un probl\`eme d'optimisation quadratique sous contraintes. Markowitz recevra le prix Nobel d'\'economie en 1990 pour ce travail. Ce chapitre d\'eveloppe la th\'eorie moyenne-variance, la fronti\`ere efficiente, le mod\`ele CAPM de Sharpe, et les extensions modernes incluant le risque CVaR et l'optimisation robuste.

\section{Th\'eorie de Markowitz}

\subsection{Cadre moyenne-variance}

\begin{definition}[Portefeuille]
Consid\'erons $n$ actifs de rendements al\'eatoires $R = (R_1, \ldots, R_n)^\top$
avec vecteur de rendements esp\'er\'es $\mu = \E[R] \in \R^n$ et matrice
de covariance $\Sigma = \Cov(R) \in \R^{n \times n}$ (sym\'etrique d\'efinie
positive).

Un \textbf{portefeuille} est d\'efini par un vecteur de poids
$w = (w_1, \ldots, w_n)^\top$ avec $\mathbf{1}^\top w = 1$. Le rendement
du portefeuille est~:
\[
    R_p = w^\top R, \quad \E[R_p] = w^\top\mu, \quad
    \Var(R_p) = w^\top\Sigma w.
\]
\end{definition}

\subsection{Probl\`eme d'optimisation}

\begin{definition}[Probl\`eme de Markowitz]
Le \textbf{probl\`eme de Markowitz} (1952) cherche le portefeuille de
variance minimale pour un rendement cible $\mu_p$~:
\[
    \min_w \frac{1}{2}w^\top\Sigma w \quad \text{s.c.} \quad
    w^\top\mu = \mu_p, \quad \mathbf{1}^\top w = 1.
\]
\end{definition}

\begin{theoreme}[Solution analytique]
\label{thm:markowitz_solution}
La solution du probl\`eme de Markowitz est~:
\[
    w^* = \Sigma^{-1}\bigl[\lambda_1\mu + \lambda_2\mathbf{1}\bigr],
\]
o\`u les multiplicateurs de Lagrange $\lambda_1, \lambda_2$ sont
d\'etermin\'es par les contraintes. En posant~:
\begin{align}
    A &= \mathbf{1}^\top\Sigma^{-1}\mu, \quad
    B = \mu^\top\Sigma^{-1}\mu, \quad
    C = \mathbf{1}^\top\Sigma^{-1}\mathbf{1}, \quad
    D = BC - A^2,
\end{align}
on obtient~:
\[
    \lambda_1 = \frac{C\mu_p - A}{D}, \quad
    \lambda_2 = \frac{B - A\mu_p}{D}.
\]
\end{theoreme}

\begin{formules}[title=\textbf{Variance minimale du portefeuille optimal}]
\[
    \sigma_p^2 = \frac{C\mu_p^2 - 2A\mu_p + B}{D}.
\]
C'est une parabole en $\mu_p$ dans le plan $(\sigma_p^2, \mu_p)$,
ou une hyperbole dans le plan $(\sigma_p, \mu_p)$.
\end{formules}

\begin{definition}[Fronti\`ere efficiente]
La \textbf{fronti\`ere efficiente} est la partie sup\'erieure de
l'hyperbole~: l'ensemble des portefeuilles optimaux pour lesquels il est
impossible d'augmenter le rendement sans augmenter le risque.
\end{definition}

\begin{definition}[Portefeuille de variance minimale globale (GMV)]
Le portefeuille \textbf{GMV} minimise la variance sans contrainte de
rendement~:
\[
    w_{\text{GMV}} = \frac{\Sigma^{-1}\mathbf{1}}
    {\mathbf{1}^\top\Sigma^{-1}\mathbf{1}}.
\]
Son rendement esp\'er\'e est $\mu_{\text{GMV}} = A/C$ et sa variance
$\sigma_{\text{GMV}}^2 = 1/C$.
\end{definition}

\begin{codeblock}[title=\textbf{Fronti\`ere efficiente de Markowitz}]
\begin{minted}{python}
import numpy as np
from scipy.optimize import minimize

def efficient_frontier(mu, Sigma, n_points=100):
    """Calcul de la frontiere efficiente."""
    n = len(mu)
    # Portefeuille GMV
    Sigma_inv = np.linalg.inv(Sigma)
    ones = np.ones(n)
    w_gmv = Sigma_inv @ ones / (ones @ Sigma_inv @ ones)
    mu_gmv = w_gmv @ mu

    # Plage de rendements cibles
    mu_range = np.linspace(mu_gmv - 0.02, max(mu) + 0.05, n_points)
    sigmas = []
    weights = []

    for mu_target in mu_range:
        constraints = [
            {'type': 'eq', 'fun': lambda w: np.sum(w) - 1},
            {'type': 'eq', 'fun': lambda w, m=mu_target: w @ mu - m}
        ]
        result = minimize(
            lambda w: 0.5 * w @ Sigma @ w,
            np.ones(n)/n, method='SLSQP',
            constraints=constraints
        )
        sigmas.append(np.sqrt(result.fun * 2))
        weights.append(result.x)

    return mu_range, np.array(sigmas), np.array(weights)

# Exemple avec 3 actifs
mu = np.array([0.08, 0.12, 0.15])
Sigma = np.array([[0.04, 0.006, 0.002],
                   [0.006, 0.09, 0.009],
                   [0.002, 0.009, 0.16]])

mu_range, sigmas, weights = efficient_frontier(mu, Sigma)
print(f"GMV: sigma={sigmas.min():.4f}")
\end{minted}
\end{codeblock}

\section{Mod\`ele CAPM}

\begin{theoreme}[Capital Asset Pricing Model (Sharpe, 1964)]
\`A l'\'equilibre, le rendement esp\'er\'e de tout actif $i$ satisfait~:
\[
    \E[R_i] - r = \beta_i(\E[R_M] - r),
\]
o\`u $\beta_i = \Cov(R_i, R_M)/\Var(R_M)$ est le \textbf{b\^eta} de l'actif
par rapport au portefeuille de march\'e $M$, et $r$ le taux sans risque.
\end{theoreme}

\begin{definition}[Ratio de Sharpe]
Le \textbf{ratio de Sharpe} d'un portefeuille $p$ est~:
\[
    \text{SR}_p = \frac{\E[R_p] - r}{\sigma_p}.
\]
Le portefeuille tangent (intersection de la CML avec la fronti\`ere
efficiente) maximise le ratio de Sharpe.
\end{definition}

\begin{codeblock}[title=\textbf{Portefeuille tangent (max Sharpe)}]
\begin{minted}{python}
def tangent_portfolio(mu, Sigma, rf):
    """Portefeuille tangent maximisant le ratio de Sharpe."""
    n = len(mu)
    Sigma_inv = np.linalg.inv(Sigma)
    excess = mu - rf
    w_tan = Sigma_inv @ excess / (np.ones(n) @ Sigma_inv @ excess)
    mu_tan = w_tan @ mu
    sigma_tan = np.sqrt(w_tan @ Sigma @ w_tan)
    sharpe = (mu_tan - rf) / sigma_tan
    return w_tan, mu_tan, sigma_tan, sharpe

rf = 0.03
w_tan, mu_tan, sig_tan, sr = tangent_portfolio(mu, Sigma, rf)
print(f"Portefeuille tangent: w = {w_tan}")
print(f"Rendement: {mu_tan:.4f}, Risque: {sig_tan:.4f}")
print(f"Ratio de Sharpe: {sr:.4f}")
\end{minted}
\end{codeblock}

\section{Mod\`ele de Black--Litterman}

\begin{definition}[Mod\`ele de Black--Litterman (1992)]
Le mod\`ele combine les rendements d'\'equilibre (CAPM) $\Pi = \delta\Sigma w_M$
avec les \textbf{vues} de l'investisseur $P\mu = Q + \epsilon$,
$\epsilon \sim \mathcal{N}(0, \Omega)$, pour obtenir des rendements
post\'erieurs~:
\[
    \hat{\mu} = \bigl[(\tau\Sigma)^{-1} + P^\top\Omega^{-1}P\bigr]^{-1}
    \bigl[(\tau\Sigma)^{-1}\Pi + P^\top\Omega^{-1}Q\bigr].
\]
Le portefeuille optimal utilise alors $\hat{\mu}$ dans l'optimisation
de Markowitz.
\end{definition}

\begin{intuition}[title=\textbf{Avantage de Black--Litterman}]
Black--Litterman r\'esout deux probl\`emes de Markowitz classique~:
(1)~les poids extr\^emes dus \`a la sensibilit\'e aux rendements
estim\'es, et (2)~l'absence de lien avec l'\'equilibre de march\'e.
Les vues sont exprim\'ees de mani\`ere intuitive et le mod\`ele
g\'en\`ere des portefeuilles plus stables.
\end{intuition}

\section{Contraintes pratiques}

\begin{remarque}
En pratique, on ajoute souvent des contraintes suppl\'ementaires~:
\begin{itemize}
    \item Interdiction de la vente \`a d\'ecouvert~: $w_i \geq 0$.
    \item Limites de concentration~: $w_i \leq w_{\max}$.
    \item Contraintes sectorielles ou de liquidit\'e.
    \item Co\^uts de transaction (optimisation dynamique).
\end{itemize}
\end{remarque}

\begin{codeblock}[title=\textbf{Optimisation avec contraintes}]
\begin{minted}{python}
def efficient_frontier_constrained(mu, Sigma, rf, n_points=50):
    """Frontiere efficiente avec w >= 0 (pas de vente a decouvert)."""
    n = len(mu)
    bounds = [(0, 1) for _ in range(n)]
    mu_range = np.linspace(min(mu), max(mu), n_points)
    results = []

    for mu_target in mu_range:
        constraints = [
            {'type': 'eq', 'fun': lambda w: np.sum(w) - 1},
            {'type': 'eq', 'fun': lambda w, m=mu_target: w @ mu - m}
        ]
        res = minimize(lambda w: w @ Sigma @ w, np.ones(n)/n,
                       method='SLSQP', bounds=bounds,
                       constraints=constraints)
        if res.success:
            results.append((mu_target, np.sqrt(res.fun), res.x))

    return results
\end{minted}
\end{codeblock}

\section{Optimisation robuste}

\begin{definition}[Optimisation robuste]
L'\textbf{optimisation robuste} tient compte de l'incertitude sur les
param\`etres estim\'es. On r\'esout~:
\[
    \min_w \max_{\mu \in \mathcal{U}} \left\{-w^\top\mu
    + \lambda\,w^\top\Sigma w\right\},
\]
o\`u $\mathcal{U}$ est un ensemble d'incertitude autour de $\hat{\mu}$.
Pour $\mathcal{U} = \{\mu : \norm{\mu - \hat{\mu}}_{\Sigma^{-1}} \leq \kappa\}$,
le probl\`eme se r\'eduit \`a~:
\[
    \min_w \left\{-w^\top\hat{\mu} + \kappa\sqrt{w^\top\Sigma w}
    + \lambda\,w^\top\Sigma w\right\}.
\]
\end{definition}

\section{Estimation des param\`etres}

\begin{attention}[title=\textbf{Erreurs d'estimation}]
Les rendements esp\'er\'es $\mu$ sont notoirement difficiles \`a estimer.
Avec $T$ observations, l'erreur standard de $\hat{\mu}_i$ est
$\sigma_i/\sqrt{T}$, souvent du m\^eme ordre que le signal.
La matrice $\Sigma$ est mieux estim\'ee mais peut \^etre singuli\`ere
ou mal conditionn\'ee lorsque $n \approx T$.
\end{attention}

\begin{proposition}[Estimateur de Ledoit--Wolf]
L'estimateur \textbf{shrinkage} de Ledoit--Wolf (2004) r\'egularise la
matrice de covariance~:
\[
    \hat{\Sigma}_{\text{LW}} = (1-\alpha)\hat{\Sigma} + \alpha\,\nu I,
\]
o\`u $\alpha \in [0,1]$ et $\nu = \text{tr}(\hat{\Sigma})/n$ sont
d\'etermin\'es analytiquement pour minimiser la perte quadratique.
\end{proposition}

\section{Exercices}

\begin{exercice}[Fronti\`ere efficiente]
\begin{enumerate}
    \item Calculer et tracer la fronti\`ere efficiente pour 5 actifs.
    \item Identifier le portefeuille GMV et le portefeuille tangent.
    \item Tracer la Capital Market Line.
    \item Comparer les fronti\`eres avec et sans contrainte de positivit\'e.
\end{enumerate}
\end{exercice}

\begin{exercice}[CAPM et b\^eta]
\begin{enumerate}
    \item Estimer les b\^etas de 5 actions par r\'egression lin\'eaire sur
          un indice de march\'e.
    \item V\'erifier la relation du CAPM en tra\c{c}ant le rendement moyen
          historique en fonction du b\^eta.
    \item Calculer l'alpha de Jensen pour chaque action.
\end{enumerate}
\end{exercice}

\begin{exercice}[Black--Litterman]
\begin{enumerate}
    \item Impl\'ementer le mod\`ele de Black--Litterman.
    \item Partir des rendements d'\'equilibre et ajouter la vue ``l'actif 1
          surperformera l'actif 2 de 2\%''.
    \item Comparer les poids avant et apr\`es l'int\'egration de la vue.
\end{enumerate}
\end{exercice}

\begin{exercice}[Estimation robuste]
\begin{enumerate}
    \item Simuler 250 jours de rendements pour 10 actifs et estimer
          $\hat{\mu}$ et $\hat{\Sigma}$.
    \item Comparer les portefeuilles optimaux avec $\hat{\Sigma}$ et
          $\hat{\Sigma}_{\text{LW}}$.
    \item Effectuer un backtest et mesurer la performance out-of-sample.
\end{enumerate}
\end{exercice}
